
\documentclass[conference]{IEEEtran}

\usepackage[utf8]{inputenc}
\usepackage[T1]{fontenc}
\usepackage{cite}
\usepackage{booktabs}
\usepackage{array}
\usepackage{url}
\usepackage{hyperref}
\usepackage{microtype}
\usepackage{amsmath}
\usepackage{enumitem}
\usepackage{xcolor}

\hypersetup{
  colorlinks=true,
  linkcolor=black,
  citecolor=black,
  urlcolor=blue
}

\newcommand{\bert}{\texttt{BERT}}
\newcommand{\lr}{\mathrm{LR}}

\begin{document}

\title{Adapting BERT for Four Classical NLP Tasks:\\
An Empirical Comparison of Feature-Based, Partial, and Full Fine-Tuning}

\author{
\IEEEauthorblockN{
Dami\'an Nicol\'as S\'anchez Novelo,
Russel Emmanuel Ku Aguilar,
Jonathan Abisai Velasco Mart\'in,\\
\'Angel Rivaldo Canch\'e Puc,
Bianca Alexandra Acosta Castellanos
}
\IEEEauthorblockA{
Data Engineering, Universidad Polit\'ecnica de Yucat\'an\\
Ucu, Yucat\'an, M\'exico\\
U2T01 -- Adapting BERT for NLP Tasks
}
}

\maketitle

\begin{abstract}
This work evaluates how different BERT adaptation strategies behave across four classical natural language processing tasks: topic classification, named entity recognition (NER), part-of-speech (POS) tagging, and extractive question answering (QA). The experiments compare feature-based adaptation, partial fine-tuning, and full fine-tuning using BERT-base models and task-specific output heads. For fine-tuning runs, the newly initialized task head and pretrained encoder are optimized with separate learning-rate groups. Results show that full fine-tuning achieved the highest measured predictive performance in all four task experiments, although the computational cost and training configuration varied substantially by task. The selected models were documented and published with tokenizers and model cards on the Hugging Face Hub.
\end{abstract}

\begin{IEEEkeywords}
BERT, transfer learning, fine-tuning, natural language processing, named entity recognition, part-of-speech tagging, question answering, topic classification
\end{IEEEkeywords}

\section{Introduction}
BERT made it practical to reuse a single pretrained Transformer encoder across many downstream NLP tasks by replacing or extending only the task-specific prediction head \cite{devlin2019bert}. This project studies three adaptation levels: feature-based adaptation, in which BERT remains frozen and its hidden representations are consumed by another learner; partial fine-tuning, in which the head and only upper encoder layers are updated; and full fine-tuning, in which the entire encoder and head are trained.

The experimental objective was not only to obtain strong scores, but also to measure what additional predictive performance is obtained as more pretrained parameters are made trainable. Four tasks were selected because they exercise different output structures: one label per document, one label per token, syntactic token labeling, and answer-span prediction. The tasks are Topic Classification on AG News, NER on CoNLL-2003, POS tagging on UD English EWT, and Extractive QA on SQuAD v1.1.

\section{Experimental Design}

\subsection{Common Principles}
All delivered neural models use a BERT-base body. Fixed random seeds were used to improve reproducibility. For runs that jointly train a newly initialized head and pretrained encoder layers, separate parameter groups are used so that the head can learn with a substantially larger learning rate than the pretrained encoder. This follows the project requirement that the randomly initialized task head should receive larger updates while the pretrained representation is protected from destructive steps.

For NER and POS, WordPiece subword alignment assigns the gold label to the first subword and uses the ignore index \texttt{-100} for continuation pieces, special tokens, and padding. Sanity checks were executed before training to verify token-label alignment. In QA, answer spans are mapped to \texttt{start\_positions} and \texttt{end\_positions}; a pre-training sanity check decoded the selected span to confirm correct alignment.

\subsection{Evaluation}
Each task compares at least two adaptation methods trained by the team. Metrics were selected according to the output structure:
\begin{itemize}[leftmargin=*]
    \item Topic Classification: accuracy and macro F1.
    \item NER: entity-level precision, recall, F1, and token accuracy.
    \item POS: token accuracy and macro precision, recall, and F1.
    \item QA: Exact Match (EM) and token-overlap F1 using the standard SQuAD normalization.
\end{itemize}
Training diagnostics include loss trajectories and measured training time. A single seeded run was used per reported configuration; therefore small gaps should not be interpreted as universal superiority independent of seed and configuration.

\section{Task-Specific Methods}

\subsection{Topic Classification}
AG News contains four balanced categories: World, Sports, Business, and Sci/Tech. The original 120,000 training examples were split with a fixed seed into 108,000 training and 12,000 validation examples; the official 7,600-example test split was reserved for final evaluation.

The first method freezes \texttt{bert-base-uncased}, extracts the 768-dimensional \texttt{[CLS]} representation, and trains Logistic Regression. The second method fully fine-tunes BERT with a four-class classification head. The encoder learning rate is $2\times10^{-5}$ and the classification-head learning rate is $10^{-3}$. The full model was trained for one epoch with a training batch size of 16, evaluation batch size of 32, and seed 42.

\subsection{Named Entity Recognition}
NER uses CoNLL-2003 \cite{tjong2003conll} and \texttt{bert-base-cased}. Partial fine-tuning trains the classification head and the top two encoder layers, while full fine-tuning updates the complete encoder and head. The corrected configuration uses separate learning rates of approximately $10^{-3}$ for the classification head and $2\times10^{-5}$ for pretrained encoder layers. Full fine-tuning was selected using validation F1 before final test evaluation.

\subsection{Part-of-Speech Tagging}
POS tagging uses the English Web Treebank configuration of Universal Dependencies \cite{nivre2020ud}. The model is \texttt{bert-base-cased} with 17 UPOS labels. The corrected partial method trains the classifier and encoder layers 10--11 with two parameter groups: $10^{-3}$ for the classifier and $2\times10^{-5}$ for the pretrained layers. The full method trains all BERT layers plus the classifier using the same two-rate principle, cosine scheduling, and warmup.

\subsection{Extractive Question Answering}
QA uses SQuAD v1.1 \cite{rajpurkar2016squad} and \texttt{bert-base-cased}. The partial method trains the QA head and the top four encoder layers; the full method trains the complete model. Both use separate parameter groups with $3\times10^{-4}$ for the QA head and $3\times10^{-5}$ for the encoder, AdamW, cosine scheduling, warmup, gradient clipping, and mixed precision when CUDA is available.

The final experiment used the full SQuAD v1.1 training split (87,599 examples) rather than the approximately 15k contingency subsample mentioned in the assignment guidance. This deviation was documented because the available hardware and mixed-precision configuration made full-dataset training practical, while also providing a direct comparison with the standard SQuAD development benchmark.

\section{Results}

\subsection{Cross-Task Comparison}
Table~\ref{tab:mainresults} summarizes the two measured approaches for each task. Because metrics differ by task, values should be interpreted within each task rather than compared numerically across tasks.

\begin{table*}[t]
\centering
\caption{Measured comparison of adaptation methods by task.}
\label{tab:mainresults}
\small
\begin{tabular}{@{}llllrrl@{}}
\toprule
Task & Method & Trainable Params. & Primary Result & Secondary Result & Time (s) & Selected \\
\midrule
Topic & Frozen BERT + Logistic Regression & Frozen encoder & Acc. 90.36\% & Macro F1 $\approx$ 0.90 & 222.54 & No \\
Topic & Full Fine-Tuning & Full model & Acc. 94.45\% & Macro F1 0.9444 & 2793.66 & Yes \\
NER & Partial FT (top 2 + head) & 14.18M & F1 88.79\% & P/R 89.03/88.54\% & 115.6 & No \\
NER & Full Fine-Tuning & 107.73M & F1 91.54\% & P/R 91.55/91.52\% & 358.5 & Yes \\
POS & Partial FT (top 2 + head) & 14.19M & Macro F1 0.8909 & Acc. 95.43\% & 173.45 & No \\
POS & Full Fine-Tuning & 107.73M & Macro F1 0.9446 & Acc. 97.42\% & 123.40 & Yes \\
QA & Partial FT (top 4 + head) & 28.35M & F1 85.17\% & EM 76.75\% & 1976.1 & No \\
QA & Full Fine-Tuning & 107.72M & F1 88.21\% & EM 80.99\% & 4075.4 & Yes \\
\bottomrule
\end{tabular}
\end{table*}

\subsection{Topic Classification}
The feature-based model achieved 90.36\% accuracy. Full fine-tuning increased accuracy to 94.45\% and macro F1 to 0.9444. The full run recorded training loss 0.180748, validation loss 0.173437, and official-test loss 0.175384.

Class-level F1 for the final model was 0.95 for World, 0.99 for Sports, 0.92 for Business, and 0.92 for Sci/Tech. The final model produced 422 errors among 7,600 official test examples (5.55\%). The dominant bidirectional confusion was Business$\leftrightarrow$Sci/Tech, including 135 Business$\rightarrow$Sci/Tech and 86 Sci/Tech$\rightarrow$Business errors. These counts are consistent with overlapping language between corporate and technology reporting, but the confusion matrix by itself does not establish a causal explanation.

\subsection{Named Entity Recognition}
Full fine-tuning improved test F1 from 88.79\% to 91.54\%, while training time increased from approximately 115.6 to 358.5 seconds and peak VRAM increased from approximately 0.76 to 2.42 GiB. For the selected model, entity-level F1 was strongest for PER (96.33\%) and LOC (93.25\%), followed by ORG (89.77\%) and MISC (80.71\%). The heterogeneous MISC class remained the most difficult category.

\subsection{Part-of-Speech Tagging}
The corrected partial model achieved 95.43\% accuracy and macro F1 0.8909. Full fine-tuning reached 97.42\% accuracy and macro F1 0.9446. The measured total times were 173.45 seconds for the partial configuration and 123.40 seconds for the full configuration. These timing results should not be treated as a causal efficiency comparison because the two configurations differ in epoch count, scheduler behavior, and training setup.

\subsection{Extractive Question Answering}
Partial fine-tuning achieved F1 85.17\% and EM 76.75\%. The selected full model achieved F1 88.21\%; the reported benchmark EM reached 80.99\%. The corresponding training times were approximately 1976.1 and 4075.4 seconds. Full fine-tuning therefore provided the stronger answer-span performance at substantially greater compute cost.

\section{Discussion}
Across the four experiments, full fine-tuning produced the highest measured predictive performance. The magnitude of the gain, however, was task dependent. Topic Classification gained approximately 4.09 percentage points in accuracy over frozen BERT features with Logistic Regression, while NER gained approximately 2.75 F1 points over partial fine-tuning. QA also improved both EM and F1, but required roughly twice the measured training time of the partial configuration.

The results illustrate why adaptation should be evaluated as a trade-off rather than as a single-score competition. Partial fine-tuning can sharply reduce the number of trainable parameters and memory use, as shown clearly in NER, while retaining much of the final performance. Feature-based adaptation was computationally simpler for Topic Classification and still achieved above 90\% accuracy. Full fine-tuning was selected for all four final deployments because it achieved the strongest task metric in each team's measured configuration.

\section{Reproducibility and Published Artifacts}
The codebase includes notebooks, configuration and requirements files, fixed seeds, training diagnostics, and documentation for reproducing the experiments. The final selected models and tokenizers were published on Hugging Face with model cards documenting training data, evaluation metrics, intended use, limitations, and references.

\begin{itemize}[leftmargin=*]
    \item NER: \url{https://huggingface.co/jonav/bert-base-cased-ner-conll2003}
    \item POS: \url{https://huggingface.co/Rivaldo2309030/bert-base-cased-pos-tagging-ewt}
    \item QA: \url{https://huggingface.co/RusselKuAguilar/bert-base-cased-squad-extractive-qa}
    \item Topic: \url{https://huggingface.co/bialexacosta21/bert-agnews-topic-classification}
    \item Project repository: \url{https://github.com/RusselKu/Bert_NLP_TaskAdapt}
\end{itemize}

\section{Limitations}
Several limitations affect interpretation. First, each reported configuration uses one fixed seed, and small score differences may fall within ordinary seed variance. Second, training configurations are not identical across tasks or even across all methods within a task; timing should therefore be interpreted descriptively. Third, Topic Classification truncates inputs to 128 tokens and uses only one fine-tuning epoch. Fourth, the QA experiment used the full SQuAD v1.1 training set rather than the suggested approximately 15k subsample. Finally, all four selected datasets and models are English-language resources, so the reported results should not be generalized to multilingual settings without additional evaluation.

\section{Conclusion}
This project demonstrates empirically how BERT adaptation depth changes the balance between performance and computational cost across four NLP tasks. Feature-based adaptation provided a strong lightweight baseline for topic classification, and partial fine-tuning substantially reduced trainable parameters for token-level and QA tasks. Nevertheless, full fine-tuning achieved the highest measured task performance in all four experiments and was therefore selected for final publication. The work also emphasizes reproducibility, correct subword handling, task-appropriate metrics, explicit training diagnostics, and model documentation as necessary components of credible transfer-learning experiments.

\begin{thebibliography}{00}

\bibitem{devlin2019bert}
J. Devlin, M.-W. Chang, K. Lee, and K. Toutanova,
``BERT: Pre-training of Deep Bidirectional Transformers for Language Understanding,''
in \emph{Proc. NAACL-HLT}, 2019, pp. 4171--4186.

\bibitem{tjong2003conll}
E. F. Tjong Kim Sang and F. De Meulder,
``Introduction to the CoNLL-2003 Shared Task: Language-Independent Named Entity Recognition,''
in \emph{Proc. CoNLL}, 2003, pp. 142--147.

\bibitem{nivre2020ud}
J. Nivre \emph{et al.},
``Universal Dependencies v2: An Evergrowing Multilingual Treebank Collection,''
in \emph{Proc. LREC}, 2020, pp. 4034--4043.

\bibitem{rajpurkar2016squad}
P. Rajpurkar, J. Zhang, K. Lopyrev, and P. Liang,
``SQuAD: 100,000+ Questions for Machine Comprehension of Text,''
in \emph{Proc. EMNLP}, 2016, pp. 2383--2392.

\bibitem{zhang2015char}
X. Zhang, J. Zhao, and Y. LeCun,
``Character-level Convolutional Networks for Text Classification,''
in \emph{Advances in Neural Information Processing Systems}, vol. 28, 2015.

\bibitem{wolf2020transformers}
T. Wolf \emph{et al.},
``Transformers: State-of-the-Art Natural Language Processing,''
in \emph{Proc. EMNLP: System Demonstrations}, 2020, pp. 38--45.

\end{thebibliography}

\end{document}
